\documentclass[11pt]{article}

\usepackage[margin=1in]{geometry}
\usepackage{amsmath,amssymb,amsthm,mathtools}
\usepackage{booktabs}
\usepackage{microtype}
\usepackage{xcolor}
\usepackage{hyperref}
\usepackage{enumitem}

\hypersetup{
  colorlinks=true,
  linkcolor=blue!55!black,
  citecolor=blue!55!black,
  urlcolor=blue!55!black
}

\newtheorem{theorem}{Theorem}[section]
\newtheorem{proposition}[theorem]{Proposition}
\newtheorem{corollary}[theorem]{Corollary}
\newtheorem{lemma}[theorem]{Lemma}
\theoremstyle{definition}
\newtheorem{definition}[theorem]{Definition}
\theoremstyle{remark}
\newtheorem{remark}[theorem]{Remark}

\newcommand{\br}[2]{[#1#2]}
\newcommand{\abs}[1]{\lvert #1\rvert}
\newcommand{\Vbar}{\overline V}
\newcommand{\Mhat}{\widehat{\mathrm T}}
\newcommand{\R}{\mathbb{R}}

\title{Half-Collinear Graviton Vertices as\\Weighted Root-Cone Enumerators}
\author{James Kehoe\\
\small Independent researcher}
\date{August 9, 2026}

\begin{document}
\maketitle

\begin{abstract}
The general half-collinear single-minus graviton amplitude recently obtained by
Guevara, Lupsasca, Skinner, Strominger, and Weil is expressed through retarded
and advanced multipoint vertices whose Cayley-tree weights contain nonlocal
step functions.  We show that each retarded vertex is equivalently a weighted
enumerator of directed spanning-tree root cones containing a kinematically
defined netflow vector; the advanced vertex evaluates the opposite cones.
Consequently, the kinematic walls are contained among the hyperplanes of the
resonance arrangement.  This resolves the published matrix-tree obstruction
in a class of chambers: the global cut conditions become root-cone membership,
and whenever the feasible family is arborescent the ordinary directed
matrix-tree theorem applies.  After ordering the projective slopes of the
block spinors, we characterize which labeled resonance chambers are realizable
by generic real two-spinors and show that their number is
$(m-1)R_{m-1}/m$ for an $m$-block vertex, where $R_{m-1}$ is the number of
rank-$(m-1)$ resonance chambers.  An exhaustive
classification gives 24 realizable chambers at four blocks and 296 at five
blocks.  Six nontrivial four-block chambers and 57 nontrivial five-block
chambers reduce to directed arborescence enumerators and therefore admit
matrix-tree determinant formulas.  As an explicit application, we derive a
compact five-graviton formula in a chamber outside the previously studied
decay region.  All low-rank classifications and the five-point identity are
verified with exact arithmetic.
\end{abstract}

\section{Introduction}

Reference~\cite{GuevaraGraviton} established that single-minus tree-level
graviton amplitudes, often presumed to vanish, are nonzero on half-collinear
configurations in Klein signature or at complexified momenta.  The authors
derived an on-shell Berends--Giele recursion and an explicit solution in terms
of unordered Feynman trees~\cite{BerendsGiele}.  In a special decay chamber,
the answer collapses to a product of soft factors after an application of the
directed matrix-tree theorem, in parallel with determinant formulas familiar
from graviton amplitudes~\cite{Hodges}.  Simplification in general kinematics
was left open; the source paper also notes that the global tree dependence
prevents a direct matrix-tree treatment of its general formula.

The obstruction is easy to state.  A Cayley tree contributes only when every
edge passes a sign test involving the total momentum on both sides of the cut
created by deleting that edge.  These tests depend on the global tree, so the
usual edge-local matrix-tree theorem does not immediately apply.

Our main observation is that the global cut test has a standard interpretation:
it asks whether a prescribed netflow is generated by strictly positive flows
on the directed tree.  This identifies the graviton vertex with a weighted
root-cone enumerator and imports the chamber combinatorics of the resonance
arrangement~\cite{RootCones}.  The reformulation produces three concrete gains:

\begin{enumerate}[leftmargin=2em]
  \item a geometric containment and chamber classification of the kinematic walls;
  \item a count and constructive parameterization of real-spinor-realizable
        low-rank chambers; and
  \item new determinant and compact amplitude formulas in chambers beyond the
        decay region.
\end{enumerate}

The appearance of this arrangement around objects called ``retarded'' is not
without precedent.  Resonance chambers count generalized retarded functions
in axiomatic and thermal quantum field theory~\cite{Evans,EpsteinTrees}, and
the adjoint braid arrangement supports a modern combinatorial formulation of
Steinmann relations and their factorization properties
\cite{LiuNorledgeOcneanu,NorledgeOcneanu}.  Those works do not identify the
specific half-collinear graviton vertex studied here with a weighted root-cone
enumerator.  They nevertheless show that the general link between retarded
objects and resonance chambers is classical; the new claims below are the
explicit flow-cone equivalence, the realizability theorem, the arborescent
classification, and the non-decay five-point formula.

This manuscript is a working research draft.  Its claims and computations have
passed an adversarial independent machine review, but they have not yet been
checked by an unaffiliated human expert in scattering amplitudes or resonance
arrangements.

\section{The half-collinear multipoint vertex}

Let $q_1,\ldots,q_m\in\R^2$ denote the real two-spinors entering a block
vertex, let
\begin{equation}
  Q=\sum_{i=1}^m q_i,
  \qquad [A,B]=\left[\sum_{i\in A}q_i,\sum_{j\in B}q_j\right],
\end{equation}
and write $[ij]=[q_i,q_j]$ for the antisymmetric determinant.  We assume
generic data: no pair bracket and no cut bracket appearing below vanishes.

For a Cayley tree $T$ on $[m]=\{1,\ldots,m\}$, deleting an edge
$e=\{u,v\}$ produces a cut $A_e\sqcup B_e=[m]$.  Choosing $u\in A_e$, the
retarded vertex of Ref.~\cite{GuevaraGraviton} is
\begin{equation}\label{eq:publishedV}
  V(q_1,\ldots,q_m)=
  \sum_{T}\prod_{e=\{u,v\}\in E(T)}
  \abs{[uv]}\,
  \Theta\!\left(-\frac{[A_e,B_e]}{[uv]}\right).
\end{equation}
The choice is immaterial: simultaneously exchanging
$(u,A_e)\leftrightarrow(v,B_e)$ reverses both numerator and denominator in the
step-function argument.  This agrees with the source convention, which fixes
$u$ to be the smaller endpoint label.
The advanced vertex $\Vbar$ is obtained by reversing the sign in every
step-function argument, and the LSZ kernel is $\Mhat=V-\Vbar$.

The apparent nonlocality lies entirely in $[A_e,B_e]$, which depends on the
whole component produced by cutting $e$.

\section{Root-cone representation}

Define the projected vector $x\in\R^m$ by
\begin{equation}\label{eq:xdef}
  x_i=[q_i,Q].
\end{equation}
Antisymmetry gives $\sum_i x_i=[Q,Q]=0$.  Direct every edge of the complete
graph by
\begin{equation}\label{eq:orientation}
  i\longrightarrow j \quad\Longleftrightarrow\quad [ij]<0.
\end{equation}
Every undirected Cayley tree inherits this orientation.

For a directed graph $G$ on $[m]$, its open root cone is
\begin{equation}
  C^{\circ}(G)=\left\{
    \sum_{i\to j\in E(G)} f_{ij}(e_i-e_j): f_{ij}>0
  \right\}.
\end{equation}
Thus $x\in C^{\circ}(G)$ precisely when $x$ is the divergence vector of a
strictly positive flow on $G$.

\begin{theorem}[Flow-cone form of the graviton vertex]\label{thm:rootcone}
Away from chamber walls,
\begin{equation}\label{eq:rootconesum}
  V(q_1,\ldots,q_m)=
  \sum_{\substack{T\text{ spanning tree}\\x\in C^{\circ}(T)}}
  \prod_{\{i,j\}\in E(T)}\abs{[ij]}.
\end{equation}
If every tree is given the retarded orientation~\eqref{eq:orientation}, then
\begin{equation}
  \Vbar(q_1,\ldots,q_m)=
  \sum_{\substack{T\text{ spanning tree}\\-x\in C^{\circ}(T)}}
  \prod_{\{i,j\}\in E(T)}\abs{[ij]}.
\end{equation}
\end{theorem}

\begin{proof}
For the cut induced by deleting $e$,
\begin{equation}\label{eq:cutisflow}
  [A_e,B_e]=[q_{A_e},Q-q_{A_e}]
  =[q_{A_e},Q]=\sum_{i\in A_e}x_i.
\end{equation}
A divergence vector on a directed tree fixes every edge flow uniquely.  Its
flow through $e$ has magnitude $\abs{\sum_{i\in A_e}x_i}$.  If $[uv]>0$,
the retarded orientation is $v\to u$, the edge enters $A_e$, and positivity
requires $\sum_{i\in A_e}x_i<0$.  If $[uv]<0$, the edge leaves $A_e$ and
positivity requires the opposite inequality.  Both cases are exactly
\begin{equation}
  \Theta\!\left(-\frac{[A_e,B_e]}{[uv]}\right)=1.
\end{equation}
The tree survives all step functions if and only if $x\in C^{\circ}(T)$,
and its remaining weight is the product in~\eqref{eq:rootconesum}.  Reversing
every step-function sign reverses every tree orientation, whose root cone is
$-C^{\circ}(T)$.
\end{proof}

\begin{corollary}[Resonance walls]
For fixed pair-bracket signs, every vertex wall is a resonance hyperplane
\begin{equation}
  \sum_{i\in A}x_i=0,
  \qquad \varnothing\neq A\subsetneq[m].
\end{equation}
Consequently, $V$ and $\Vbar$ are weighted tree enumerators that are constant
in tree support on every resonance chamber.
\end{corollary}

\section{Which resonance chambers are realizable?}

The resonance arrangement contains abstract sign chambers that need not occur
for real spinors with a fixed projective ordering.  The realizability
restriction is simple.

Choose an $SL(2,\R)$ frame with $Q=(1,0)$, assume every $x_i\neq0$, and write
\begin{equation}\label{eq:slopeparam}
  q_i=(x_i r_i,-x_i).
\end{equation}
After relabeling, take $r_1<\cdots<r_m$.  Momentum conservation gives
\begin{equation}\label{eq:constraints}
  \sum_i x_i=0,
  \qquad
  \sum_i x_i r_i=1,
\end{equation}
while
\begin{equation}\label{eq:pairfactor}
  [ij]=x_i x_j(r_j-r_i).
\end{equation}
Thus the pair-bracket tournament is fixed by the signs of $x_i$ once the
slopes are ordered.

\begin{theorem}[Realizable chamber criterion]\label{thm:realizable}
Let $X_k=\sum_{i=1}^k x_i$.  A generic resonance chamber---so in particular
all $x_i$ and all proper subset sums, including the prefix sums, are
nonzero---has a realization
with $r_1<\cdots<r_m$ and~\eqref{eq:constraints} if and only if
\begin{equation}
  X_k<0\quad\text{for at least one }k\in\{1,\ldots,m-1\}.
\end{equation}
Equivalently, the only excluded chambers are those in the positive root cone
aligned with the slope order.
\end{theorem}

\begin{proof}
Summation by parts yields
\begin{equation}\label{eq:summationbyparts}
  \sum_{i=1}^m x_i r_i
  =-\sum_{k=1}^{m-1}X_k(r_{k+1}-r_k).
\end{equation}
Every slope gap is positive.  Hence the right-hand side can equal $1$ only if
some $X_k<0$.  Conversely, if one prefix is negative, choose its slope gap
sufficiently large relative to all other positive gaps so that the right-hand
side is positive, and rescale all gaps to make it equal to $1$.
\end{proof}

Let $R_{m-1}$ be the number of rank-$(m-1)$ resonance chambers.  The cyclic
decomposition in Ref.~\cite{RootCones} shows that exactly $R_{m-1}/m$ lie in
the aligned positive root cone.  Theorem~\ref{thm:realizable} therefore gives
the number of labeled, slope-ordered realizable chambers,
\begin{equation}\label{eq:realizablecount}
  R^{\mathrm{real}}_{m-1}=\frac{m-1}{m}R_{m-1}.
\end{equation}

\begin{remark}
Here ``realizable'' is a kinematic statement about generic real two-spinors,
not a claim that every counted chamber lies in a fixed scattering channel.
The decay region of Ref.~\cite{GuevaraGraviton} additionally imposes energy
signs and therefore selects a smaller in/out configuration.
\end{remark}

\section{Determinantal chambers and low-rank classification}

For a fixed realizable chamber $\mathcal C$, let $\mathcal F_{\mathcal C}$ be
the family of retarded-oriented spanning trees whose open root cones contain
$x$.  Theorem~\ref{thm:rootcone} is always a weighted sum over this family.
Some chambers restore an edge-local description.

\begin{definition}
A chamber is \emph{arborescent} if there are a directed graph $D_{\mathcal C}$
and a root $\rho$ such that $\mathcal F_{\mathcal C}$ is exactly the set of all
directed spanning trees of $D_{\mathcal C}$ oriented toward $\rho$.
\end{definition}

Give every arc $i\to j$ weight $w_{ij}=\abs{[ij]}$ and define the outgoing
Laplacian
\begin{equation}
  L_{ij}=\begin{cases}
    -w_{ij},&i\to j,\ i\neq j,\\
    \sum_{i\to k}w_{ik},&i=j,\\
    0,&\text{otherwise}.
  \end{cases}
\end{equation}

\begin{proposition}[Determinantal chamber formula]\label{prop:det}
In an arborescent chamber,
\begin{equation}
  V(q_1,\ldots,q_m)=\det L^{(\rho)},
\end{equation}
where $L^{(\rho)}$ deletes row and column $\rho$.
\end{proposition}

\begin{proof}
This is the directed matrix-tree theorem combined with
Theorem~\ref{thm:rootcone} and the definition of an arborescent chamber.
\end{proof}

We exhaustively enumerated the known 32 rank-three and 370 rank-four
resonance chambers.  The realizability criterion leaves 24 four-block and 296
five-block chambers.  For each realizable chamber we constructed an exact integer
representative, enumerated all $m^{m-2}$ Cayley trees, tested positive-flow
feasibility, and checked whether the resulting family was the complete
in-arborescence family of its union digraph.  Table~\ref{tab:classification}
summarizes the result.

\begin{table}[ht]
\centering
\small
\begin{tabular}{@{}lrr@{}}
\toprule
 & $m=4$ & $m=5$\\
\midrule
All resonance chambers & 32 & 370\\
Realizable ordered-slope chambers & 24 & 296\\
Distinct feasible-tree families & 9 & 55\\
Chambers with $V=0$ & 6 & 38\\
Nonempty arborescent chambers & 10 & 57\\
Nontrivial arborescent chambers & 6 & 57\\
Largest feasible-tree family & 5 & 28\\
\bottomrule
\end{tabular}
\caption{Exhaustive low-rank chamber classification.  Distinct feasible-tree
families and every subsequent row are counted only over realizable labeled
chambers.  ``Nontrivial'' means more than one contributing spanning tree.}
\label{tab:classification}
\end{table}

All 67 reported determinant identities were independently evaluated as exact
Laplacian cofactors and compared with the original step-function sums.
At five blocks every nonempty arborescent family is already nontrivial: the
observed feasible-family sizes jump directly from zero to two.

\section{A compact non-decay five-graviton amplitude}

We now lift one arborescent four-block chamber through the five-point recursion.
Let $q_1,\ldots,q_4$ be the positive-helicity block spinors,
$Q=q_1+\cdots+q_4$, and $q_5=-Q$.  Define
\begin{equation}
  x_i=[q_i,Q]=-[i5],
  \qquad
  y_i^{S}=[q_i,q_S]\quad(i\in S).
\end{equation}
Consider the open sign chamber $\mathcal C_\star$ specified in
Table~\ref{tab:signs}.  The pair signs fix the retarded tournament; the $x$
signs fix the four-block resonance chamber; and the $y$ signs fix the four
three-block subvertices entering the recursion.

\begin{table}[ht]
\centering
\small
\begin{tabular}{@{}lll@{}}
\toprule
Data & Ordered quantities & Signs\\
\midrule
Projected singles & $(x_1,x_2,x_3,x_4)$ & $(-,+,+,-)$\\
Projected pairs & $(x_{12},x_{13},x_{14},x_{23},x_{24},x_{34})$
  & $(-,-,-,+,+,+)$\\
Pair brackets & $([12],[13],[14],[23],[24],[34])$
  & $(-,-,+,+,-,-)$\\
$S=123$ & $(y_1^{123},y_2^{123},y_3^{123})$ & $(-,+,+)$\\
$S=124$ & $(y_1^{124},y_2^{124},y_4^{124})$ & $(+,+,-)$\\
$S=134$ & $(y_1^{134},y_3^{134},y_4^{134})$ & $(-,+,-)$\\
$S=234$ & $(y_2^{234},y_3^{234},y_4^{234})$ & $(+,-,+)$\\
\bottomrule
\end{tabular}
\caption{Definition of the open five-point chamber $\mathcal C_\star$.
Here $x_{ij}=x_i+x_j$.}
\label{tab:signs}
\end{table}

Writing $w_{ij}=\abs{[ij]}$, direct root-cone inspection gives
\begin{align}
  \mathcal F(V_{1234})
    &=\bigl\{\{14,23,24\},\{14,24,34\}\bigr\},\\
  \mathcal F(\Vbar_{1234})
    &=\bigl\{\{12,13,14\},\{13,14,23\}\bigr\}.
\end{align}
Both are two-term arborescence families, so
\begin{align}\label{eq:fourvertices}
  V_{1234}&=w_{14}w_{24}(w_{23}+w_{34}),\\
  \Vbar_{1234}&=w_{13}w_{14}(w_{12}+w_{23}).
\end{align}
The triple sign vectors similarly imply
\begin{equation}\label{eq:triples}
  V_{234}=V_{123}=0,
  \qquad
  V_{134}=w_{14}w_{34},
  \qquad
  V_{124}=w_{12}w_{24}.
\end{equation}

The published five-point recursion is
\begin{equation}\label{eq:M5rec}
  M_{12345}=-\Mhat_{1234}
  -w_{15}V_{234}-w_{25}V_{134}-w_{35}V_{124}-w_{45}V_{123},
\end{equation}
with $\Mhat=V-\Vbar$.  To obtain this form directly from Eq.~(33) of
Ref.~\cite{GuevaraGraviton}, use the preamplitude base cases
$\overline M_i=1$, $\overline M_{ij}=0$, and
$\overline M_{ijk}=-V_{ijk}$.  Only the four-singleton partition and the four
partitions consisting of one singleton plus one triple then survive.
Substitution proves the following.

\begin{corollary}[Five-graviton formula outside the decay chamber]
Throughout $\mathcal C_\star$,
\begin{align}\label{eq:newM5}
  M_{12345}={}&
  w_{13}w_{14}(w_{12}+w_{23})
  -w_{14}w_{24}(w_{23}+w_{34})\notag\\
  &-w_{25}w_{14}w_{34}
  -w_{35}w_{12}w_{24}.
\end{align}
This is a genuine non-decay formula: both $V_{1234}$ and
$\Vbar_{1234}$ are nonzero in $\mathcal C_\star$.
\end{corollary}

The proof is finite and analytic: Table~\ref{tab:signs} determines every
step function in the 16 four-vertex Cayley trees and the three trees of each
triple subvertex.  Equation~\eqref{eq:newM5} is therefore constant in form on
the entire open chamber, not merely on sampled kinematics.

\section{Exact computational verification}

The accompanying Python implementation uses integer two-spinors and exact
arithmetic.  No floating tolerances enter any sign, determinant, or amplitude
comparison.  We use $[a,b]=a_0b_1-a_1b_0$, globally opposite to the component
ordering used for brackets in Ref.~\cite{GuevaraGraviton}; an $n$-point
homogeneous formula therefore acquires the expected factor $(-1)^{n-2}$ when
transcribed directly.  The following checks currently pass:

\begin{itemize}[leftmargin=2em]
  \item direct assertions of
        $M_{123}=\abs{[12]}$ and
        $M_{1234}=\tfrac12(\abs{[12]}\abs{[34]}+
        \abs{[13]}\abs{[24]}+\abs{[14]}\abs{[23]})$ on 100 exact samples;
  \item agreement with the published decay-region product formula at four and
        five points on 100 exact samples, including the bracket-convention sign;
  \item permutation invariance through five points on randomized samples;
  \item agreement of the step-function and positive-flow definitions of both
        $V$ and $\Vbar$ on 500 randomized samples through valency six;
  \item complete recovery of the known 32 and 370 resonance chambers from
        bounded integer representatives;
  \item exact matrix-tree cofactor checks in all 10 four-block and 57
        five-block nonempty arborescent chambers; and
  \item 22,102 exact integer tests of~\eqref{eq:newM5} across
        $\mathcal C_\star$, obtained by bounded grid enumeration rather than
        random sampling.
\end{itemize}

Two additional referee programs were written from scratch without importing
the bundled implementation.  They independently reconstruct the published
recursion, reproduce the decay formula, verify the flow-cone equivalence on
300 exact samples, recover Table~\ref{tab:classification}, and derive the
families in~\eqref{eq:fourvertices}--\eqref{eq:triples} from the chamber signs
alone.  These machine checks are evidence against implementation coupling;
they do not replace expert proof review.

Two tempting stronger conjectures fail.  First, the full amplitude is not the
sum over local cubic-tree vertex weights suggested by the three- and four-point
answers; it fails at generic five-point kinematics.  Second, the generic
five-point answer is not a permutation-symmetric homogeneous cubic polynomial
in the ten variables $\abs{[ij]}$: the seven graph-monomial orbit sums have
rank seven on exact data, while adjoining the amplitude raises the rank to
eight.  These failures support the conclusion that resonance cut data is
essential outside special chambers.

\section{Discussion}

The root-cone formulation does not produce a single determinant valid in every
chamber.  It instead identifies the precise combinatorial object that replaces
one: a weighted collection of tree cones containing a netflow vector.  This is
useful because resonance chambers, root-cone intersections, compatible tree
collections, and Kostant chambers already have a substantial mathematical
theory~\cite{RootCones}.

Accordingly, the result should be read as a precise partial resolution of the
matrix-tree obstruction stated in Ref.~\cite{GuevaraGraviton}, not as a closed
form for the unrestricted amplitude.  The global cut tests are converted into
root-cone membership for every chamber; the ordinary matrix-tree theorem then
applies exactly when the feasible family is arborescent, a condition classified
here through five blocks.

Several next questions are concrete.

\begin{enumerate}[leftmargin=2em]
  \item Characterize arborescent chambers without exhaustive tree enumeration.
  \item Determine whether non-arborescent families admit decompositions into a
        small signed sum of Laplacian cofactors.
  \item Lift the 57 five-block determinant identities to compact six-graviton
        formulas, tracking the required lower-valency subchambers.
  \item Relate chamber adjacency to wall-crossing identities of the amplitude.
  \item Investigate whether the root-cone language clarifies the extension of
        the $\mathcal Lw_{1+\infty}$ recursion outside the decay region.
\end{enumerate}

The most important immediate task is independent verification by experts in
scattering amplitudes and resonance arrangements.  In particular, absence of
the present connection from the cited graviton paper does not by itself
establish novelty across the full literature.

\section*{Reproducibility statement}

The release is pinned to Python 3.12.3 and is accompanied by three
exact-arithmetic programs: one implementing and anchoring the published
recursion and flow-cone equivalence, one enumerating and classifying low-rank
chambers, and one exhaustively testing the five-point identity.  It also
contains the two independently written referee programs, machine-readable
certificates for every realizable four- and five-block chamber, cryptographic
checksums, and continuous-integration tests.

The conceptual cut/flow/root-cone argument and both retarded sign cases are
also checked in Lean~4.  Every principal declaration is guarded against
\texttt{sorryAx}; continuous integration rebuilds and rechecks the compiled
environment.  The manuscript correspondence is unaudited, and the
computational claims above remain outside Lean.

\section*{Author and AI contribution statement}

James Kehoe selected and directed the research problem, made the decision to
pursue and disseminate the result, and commissioned adversarial review.
OpenAI Codex generated the central root-cone reformulation, proofs, search and
verification code, and initial manuscript under that direction.  An independent
large-language-model review then re-derived the main claims using fresh code;
its report and programs are included in the repository.  No AI system is listed
as an author.  The named human author accepts responsibility for the released
manuscript and will update or withdraw claims if subsequent expert review finds
an error.

\begin{thebibliography}{9}

\bibitem{GuevaraGraviton}
A.~Guevara, A.~Lupsasca, D.~Skinner, A.~Strominger, and K.~Weil,
``Single-minus graviton tree amplitudes are nonzero,''
\href{https://arxiv.org/abs/2603.04330}{arXiv:2603.04330 [hep-th]} (2026).

\bibitem{RootCones}
S.~C. Gutekunst, K.~M\'esz\'aros, and T.~K. Petersen,
``Root cones and the resonance arrangement,''
\href{https://arxiv.org/abs/1903.06595}{arXiv:1903.06595 [math.CO]} (2019),
The Electronic Journal of Combinatorics \textbf{28} (2021), P1.12.

\bibitem{BerendsGiele}
F.~A. Berends and W.~T. Giele,
``Recursive calculations for processes with $n$ gluons,''
Nuclear Physics B \textbf{306} (1988), 759--808.

\bibitem{Hodges}
A.~Hodges,
``New expressions for gravitational scattering amplitudes,''
Journal of High Energy Physics \textbf{07} (2013), 075,
\href{https://arxiv.org/abs/1108.2227}{arXiv:1108.2227 [hep-th]}.

\bibitem{Evans}
T.~S. Evans,
``N-point finite-temperature expectation values at real times,''
Nuclear Physics B \textbf{374} (1992), 340--370.

\bibitem{EpsteinTrees}
H.~Epstein,
``Trees,''
\href{https://arxiv.org/abs/1602.00936}{arXiv:1602.00936 [math-ph]} (2016).

\bibitem{LiuNorledgeOcneanu}
Z.~Liu, W.~Norledge, and A.~Ocneanu,
``The adjoint braid arrangement as a combinatorial Lie algebra via the
Steinmann relations,''
\href{https://arxiv.org/abs/1901.03243}{arXiv:1901.03243 [math.CO]} (2019).

\bibitem{NorledgeOcneanu}
W.~Norledge and A.~Ocneanu,
``Hopf monoids, permutohedral cones, and generalized retarded functions,''
\href{https://arxiv.org/abs/1911.11736}{arXiv:1911.11736 [math.CO]} (2019).

\end{thebibliography}

\end{document}
